\subsection{Everyday embodied monitoring and adaptive support}
\label{subsec:everyday_monitoring}

Beyond episodic encounters and centralized laboratories, everyday monitoring relocates evidence acquisition to the body and the environments in which health changes unfold. The aim is not merely to collect more measurements, but to learn an individual baseline, detect meaningful departures from it and initiate proportionate verification or support \cite{Gambhir2021ContinuousMonitoring}. Wearable systems provide the most accessible implementation, combining a body interface with sensing, power, communication, signal processing and a decision unit; failures in any one component can invalidate the end-to-end inference \cite{Ates2022WearableSensors}. Consumer smartwatches have, for example, detected presymptomatic physiological deviations relative to personal baselines in COVID-19, while multimodal wearable data have supported personalized prediction of selected clinical laboratory measurements \cite{Mishra2020SmartwatchCOVID,Dunn2021WearablePredictions}. Biochemical wearables extend observation beyond motion and vital signs: an integrated electrochemical platform combined induced sweat, microfluidic handling, calibration and wireless analysis to monitor trace metabolites and nutrients during exercise and at rest \cite{Wang2022WearableBiosensor}. Implantable and ingestible devices gain access to signals that are difficult to sample reliably at the skin. A microfluidic intraocular-pressure implant enabled smartphone-based home readings across fluctuations missed by occasional clinic measurements, and active-reset affinity sensors overcame saturation to track inflammatory proteins dynamically in living animals \cite{Araci2014ImplantableIOP,Zargartalebi2024ActiveReset}. Ingestible capsules have measured intestinal gases in a human pilot, self-powered glucose sensing in the porcine small intestine and, in a first-in-human study, respiratory and cardiac activity from within the gastrointestinal tract \cite{KalantarZadeh2018IngestibleGas,DelaPaz2022IngestibleBiosensing,Traverso2023IngestibleVitals}. These studies expand the spatial, temporal and molecular reach of observation, but their translational stages differ sharply. Motion artefact, biofluid variability, sensor fouling, calibration drift, data loss, power and telemetry constraints, foreign-body effects and uncertain links between proxy signals and clinical states must be evaluated longitudinally in the intended population. Most current platforms also remain one-way monitors: an alert or risk score becomes a diagnostic loop only when it triggers an appropriate repeat measurement, confirmatory examination or escalation, and the resulting evidence updates what happens next.

Ambient intelligence moves the sensing interface from the body into the home. Passive infrared detectors, door contacts and mattress sensors can generate a ``zero-interaction digital exhaust'' describing mobility, sleep, activity patterns and functional change without requiring the user to remember, wear or charge a device \cite{Schutz2022DigitalExhaust}. Contactless radio-frequency sensing similarly enabled a model to infer Parkinson's disease and estimate severity from nocturnal breathing, including measurements acquired in the home \cite{Yang2022ParkinsonBreathing}. Connected fixtures can turn routine activities into sampling opportunities. A mountable smart-toilet prototype integrated user identification, urinalysis, uroflowmetry and stool-image analysis, illustrating how excreta could be monitored without a separate testing episode \cite{Park2020SmartToilet}; however, passive collection raises unusually persistent questions about identity, consent, data ownership and access, especially in shared households \cite{Ge2023SmartToilets}. More broadly, home monitoring may support ageing in place and chronic-condition management, but the evidence base remains heterogeneous and digital access, maintenance and workflow integration can determine who benefits \cite{Chen2023DigitalHealthAging,Facchinetti2023SmartHomes}. The home is therefore not a controlled laboratory. Visitors, pets, co-residents, altered routines, moved devices, connectivity loss and gradual sensor drift can all change the meaning or completeness of a signal. Evaluation should consequently test person-level longitudinal calibration, missingness, false-alert burden and performance across households and social contexts, as well as whether a detected change leads to timely and useful action. In many near-term systems, the operational loop will remain human-mediated: ambient sensing identifies a deviation, a patient, caregiver or clinician verifies its context, and only then is monitoring intensified or care escalated.

Adaptive support begins when observations of the person change a physical or socially situated response. In rehabilitation, wearable sensing and machine learning can estimate motor impairment, movement quality and recovery outside periodic clinic assessments, creating a state representation with which therapy can be individualized \cite{AdansDester2020PrecisionRehabilitation}. Electromechanical and robot-assisted arm training after stroke has shown modest improvements in activities of daily living, arm function and strength across randomized trials, although interventions, populations and training dose were heterogeneous and many devices delivered repetitive, predefined exercise rather than online adaptation \cite{Mehrholz2018RobotArmTraining}. Human-in-the-loop optimization demonstrates a tighter embodied cycle: assistance parameters are varied, the user's metabolic or kinematic response is measured, and the controller selects a new policy for the same person \cite{Zhang2017ExoskeletonOptimization,Ding2018SoftExosuit}. This principle has been extended from tethered laboratory systems to portable ankle exoskeletons that use wearable kinematics to personalize assistance during variable-speed outdoor walking \cite{Slade2022RealWorldExoskeleton}. The engineering result is important, but optimization in small studies of controlled walking should not be equated with durable rehabilitation benefit, fall prevention or safe adaptation across impairment, fatigue and terrain. Socially assistive and companion robots occupy a different action space, using speech, movement, reminders and structured interaction to support engagement or well-being. Meta-analyses in older adults and people living with dementia suggest feasibility and possible benefits for selected emotional or neuropsychiatric outcomes, but effects on cognition, agitation and quality of life are inconsistent and high-quality randomized evidence remains limited \cite{Pu2019SocialRobots,Yu2022DementiaSocialRobots,Hsieh2023DementiaSocialRobots}. A recent randomized external pilot in surgical wards likewise demonstrated the feasibility of deploying a predefined-dialog robot, but found limited overall effects on engagement and perceived care quality \cite{Mlakar2025SocialRobotPilot}. Across these devices, adaptation should therefore be judged against the intended functional or care outcome, not interaction frequency alone. Integrated evaluation must include state-estimation error, controller latency, physical safety, comfort and fatigue, adherence, human override, caregiver workload, privacy and inequitable fit. Everyday embodied support is thus a function-specific continuum: monitoring may remain passive, low-risk assistance may adapt locally, and consequential clinical or physical changes may require confirmation and authorization. Its defining advance is not continuous sensing or robotics in isolation, but a bounded sensing--decision--action--feedback loop that improves everyday function or care while preserving the user's agency and a safe path to human intervention.
